\metatitle{Hives and BZ-polytopes}
\metadescription{Knutson--Tao hives, Berenstein--Zelevinsky polytopes,
Littlewood--Richardson coefficients, saturation, and Ehrhart stretching.}
\metakeywords{hive polytope,Berenstein--Zelevinsky polytope,
Littlewood--Richardson coefficient,saturation theorem,
stretched Littlewood--Richardson coefficient}

\section[hivePolytopes]{Hives and BZ-polytopes}

\defin{Knutson--Tao hives}\label{hiveDefinition} and
\defin{Berenstein--Zelevinsky patterns}\label{bzPatternDefinition} are
polyhedral models for
\hyperref[schurLittlewoodRichardson]{Littlewood--Richardson coefficients}.
They replace tableaux by lattice points in a polytope whose boundary data is
given by three partitions.  The main references are
\cite{BerensteinZelevinsky1988,BerensteinZelevinsky1992,KnutsonTau1999}.

For fixed partitions \(\lambda,\mu,\nu\), the corresponding
\defin{hive polytope}\label{hivePolytope} is cut out by linear inequalities
on a triangular array.  The boundary differences encode
\(\lambda,\mu,\nu\), and each elementary rhombus imposes a concavity
inequality.

\svgimg[width=0.72\textwidth]{svg-images/hive-rhombus-inequalities.svg}{
A hive is a triangular array satisfying one rhombus inequality in each
orientation.}

Equivalently, for each elementary rhombus with opposite vertices \(a,d\)
and \(b,c\), the corresponding inequality has the form
\[
  a+d \geq b+c.
\]

With the usual boundary convention,
\[
  c^\nu_{\lambda,\mu}
  =
  |\{\text{integer hives with boundary }(\lambda,\mu,\nu)\}|.
\]
Equivalently, one may use
\hyperref[bzPatterns]{Berenstein--Zelevinsky patterns}; these are different
coordinates for the same Littlewood--Richardson cone.

\begin{example}
The product
\[
  \schurS_{(2)}\schurS_{(1)}
  =
  \schurS_{(3)}+\schurS_{(2,1)}
\]
says that the two relevant hive polytopes each contain one integer point, and
all other boundary triples for this product contain none.
\end{example}

\subsection[hiveSaturation]{Saturation}

The hive model gives a geometric proof of the saturation theorem:
\[
  c^\nu_{\lambda,\mu}>0
  \quad\Longleftrightarrow\quad
  c^{k\nu}_{k\lambda,k\mu}>0
  \qquad (k\geq 1).
\]
This was proved by \name{Allen Knutson} and \name{Terence Tao}
\cite{KnutsonTau1999}; see also \cite{Buch2000}.  In hive language, scaling
the boundary scales the polytope.  Saturation says that rational feasibility
already detects integral feasibility.

\subsection[hiveStretching]{Stretching and Ehrhart theory}

The stretching function
\[
  k\mapsto c^{k\nu}_{k\lambda,k\mu}
\]
is naturally an Ehrhart-type lattice-point count.  The polynomiality of this
function is nontrivial and was proved by \name{Harm Derksen} and
\name{Jerzy Weyman} \cite{DerksenWeyman2002}, and by
\name{Etienne Rassart} \cite{Rassart2004}.  The polynomial is called a
\defin{stretched Littlewood--Richardson polynomial}
\label{stretchedLRPolynomial}.

\begin{conjecture}[King--Tollu--Toumazet, \cite{KingTolluToumazet2004}]
The stretched Littlewood--Richardson polynomial
\[
  k\mapsto c^{k\nu}_{k\lambda,k\mu}
\]
has nonnegative coefficients.
\end{conjecture}

This is one of the motivating examples for the
\hyperref[ehrhartConjecturalPositivity]{Ehrhart positivity conjectures} on
the Ehrhart page.  The same circle of ideas also contains stretched Kostka
coefficients, flagged skew Schur coefficients, and
\hyperref[keyEhrhart]{Ehrhart polynomials for key polynomials}.

\subsection[hiveRelatedModels]{Related models}

Hives are closely related to several other models:
\begin{itemize}
\item \hyperref[lrTableaux]{Littlewood--Richardson tableaux} give the
classical tableau model for the same coefficients.
\item \hyperref[bzPatterns]{BZ-patterns} and
\hyperref[gtLR]{GT-patterns with Yamanouchi inequalities} give alternate
polyhedral coordinates.
\item \hyperref[knutsonTaoWoodwardPuzzles]{Knutson--Tao--Woodward puzzles}
give a planar combinatorial model tied to the facets of the
Littlewood--Richardson cone \cite{KnutsonTaoWoodward2004}.
\item \hyperref[kostantPartitionFunction]{Kostant partition functions} give
piecewise-polynomial formulas for the same multiplicities.
\end{itemize}
